\chapter{Basic group theory}

\section{Structure preserving maps}

\defn{Group homomorphism}{
    Let \(G\) and \(H\) be groups. A \textbf{group homomorphism} is a map
    \[
        f : G \to H
    \]
    such that
    \[
        f(g_1 \circ_G g_2) = f(g_1) \circ_H f(g_2)
        \qquad \forall g_1,g_2 \in G.
    \]
}

\cor{
    If \(f : G \to H\) is a group homomorphism, then
    \[
        f(e_G) = e_H.
    \]
}

\cor{
    If \(f : G \to H\) and \(\tilde{f} : H \to K\) are group homomorphisms, then
    \[
        \tilde{f} \circ f : G \to K
    \]
    is again a group homomorphism.
}

\cor{
    If \(f : G \to H\) is a bijective group homomorphism, then
    \[
        f^{-1} : H \to G
    \]
    is also a group homomorphism.
}

\corp{
    Every group homomorphism preserves inverses, i.e. if \(f : G \to H\) is a
    group homomorphism, then
    \[
        f(g^{-1}) = f(g)^{-1}
        \qquad \forall g \in G.
    \]
}{
    Since \(g g^{-1} = e_G\), we get
    \[
        f(g)f(g^{-1}) = f(gg^{-1}) = f(e_G) = e_H.
    \]
    Similarly,
    \[
        f(g^{-1})f(g) = f(g^{-1}g) = f(e_G) = e_H.
    \]
    Hence \(f(g^{-1})\) is the inverse of \(f(g)\), so \(f(g^{-1}) = f(g)^{-1}\).
}{}

\defn{Group isomorphism}{
    An invertible group homomorphism is called a \textbf{group isomorphism}.
    If there exists a group isomorphism
    \[
        f : G \to H
    \]
    then we say that \(G\) and \(H\) are \textbf{isomorphic}.
}

\ex{
    Examples of discrete groups:
    \begin{itemize}
        \item \((\mathbb{Z},+)\).
        \item \((\mathbb{R},+)\) and \((\mathbb{C},+)\).
        \item \((\mathbb{R}\setminus\{0\},\cdot)\) and \((\mathbb{C}\setminus\{0\},\cdot)\).
        \item Let \(n \in \mathbb{N}_{>0}\). Then \(\mathbb{Z}/n\mathbb{Z}\) is the quotient of
        \(\mathbb{Z}\) by the equivalence relation
        \[
            a \sim b
            \qquad \Longleftrightarrow \qquad
            a = b + kn
            \quad \text{for some } k \in \mathbb{Z}.
        \]
        The equivalence class of \(a\) is \( [a] = \{ b \in \mathbb{Z} \mid b \sim a \} \), and
        \(\mathbb{Z}/n\mathbb{Z}\) becomes a group with operation
        \[
            [a]\cdot[b] = [a+b].
        \]
        \item The permutation group \(S_n\), for \(n \in \mathbb{N}_{>0}\), is the
        group of all permutations of \(\{1,2,\dots,n\}\). With composition as
        multiplication, \(S_n\) forms a group, and \( |S_n| = n! \).

        For example, in \(S_3\) let
        \[
            \sigma_1 = \permu{1,2,3}{3,2,1},
            \qquad
            \sigma_2 = \permu{1,2,3}{1,3,2}
        \]
        Then
        \[
            \sigma_2 \circ \sigma_1 = \permu{1,2,3}{3,1,2}.
        \]
    \end{itemize}
}

\ex{
    Examples of continuous groups:
    \begin{itemize}
        \item The general linear group
        \[
            \mathrm{GL}(n,\mathbb{R})
            =
            \{ M : \mathbb{R}^n \to \mathbb{R}^n \mid M \text{ linear},\ \det M \neq 0 \}.
        \]
        \item The orthogonal group
        \[
            \mathrm{O}(n)
            =
            \{ M : \mathbb{R}^n \to \mathbb{R}^n \mid M \text{ linear},\ M^T M = I \}.
        \]
        Equivalently, \(M \in \mathrm{O}(n)\) iff \(\det M = \pm 1\).
        \item More generally, let \(T_{ij} \in T\mathbb{R}^{n*} \otimes T\mathbb{R}^{n*}\) be a tensor.
        Its symmetry group is
        \[
            \operatorname{Sym}(T_{ij})
            =
            \{ M^i{}_j \in \mathrm{GL}(n,\mathbb{R}) \mid T_{ij} = T_{ab} M^a{}_i M^b{}_j \}.
        \]
        If \(\langle v,w\rangle = T_{ij}v^i w^j\), this is equivalently the group preserving the
        bilinear form defined by \(T\).

        For
        \[
            T =
            \begin{pmatrix}
                1 & 0 & 0 & 0 \\
                0 & -1 & 0 & 0 \\
                0 & 0 & -1 & 0 \\
                0 & 0 & 0 & -1
            \end{pmatrix},
        \]
        one gets \(\operatorname{Sym}(T) = \mathrm{O}(1,3)\). For the standard symplectic form
        \[
            T =
            \begin{pmatrix}
                0 & I_n \\
                -I_n & 0
            \end{pmatrix},
        \]
        one gets the symplectic group.
        \item An \textbf{affine transformation} is a map \(\mathbb{R}^n \to \mathbb{R}^n\) of the form
        \[
            x \mapsto Mx + T,
            \qquad M \in \mathrm{GL}(n,\mathbb{R}),\ T \in \mathbb{R}^n.
        \]
        We denote the group of affine transformations by \(\operatorname{Aff}(\mathbb{R}^n)\).
        \item If \(B \subseteq \mathbb{R}^n\), define
        \[
            \operatorname{Sym}(B)
            =
            \{ A \in \operatorname{Aff}(\mathbb{R}^n) \mid A(B) = B \}.
        \]
        For an equilateral triangle \(\Delta\), one has \(\operatorname{Sym}(\Delta) \cong S_3\).
    \end{itemize}
}

\defn{Subgroup}{
    Let \(G\) be a group. A \textbf{subgroup} \(H \leq G\) is a subset \(H \subseteq G\)
    such that \(H\) is itself a group:
    \begin{enumerate}
        \item \(h \circ_G h' \in H\) for all \(h,h' \in H\),
        \item if \(h \in H\), then \(h^{-1} \in H\).
    \end{enumerate}
}

\ex{
    Examples:
    \begin{itemize}
        \item \(\mathrm{SO}(n) \leq \mathrm{O}(n) \leq \mathrm{GL}(n)\).
        \item \(S_2 \leq S_3 \leq S_4\).
    \end{itemize}

    This notation is slightly sloppy: strictly speaking, one should specify an
    injective group homomorphism \(H \hookrightarrow G\). For example,
    \[
        H = \{ \sigma \in S_4 \mid \sigma(4) = 4 \}
    \]
    is isomorphic to \(S_3\).
}

\propp{
    Let \(f : G \to H\) be a group homomorphism. Then the kernel of \(f\),
    \[
        \ker f := \{ g \in G \mid f(g) = e_H \} \leq G,
    \]
    is a subgroup. Likewise, the image of \(f\),
    \[
        \operatorname{im} f := \{ h \in H \mid \exists g \in G,\ f(g)=h \} \leq H,
    \]
    is a subgroup.
}{
    Let \(g_1,g_2 \in \ker f\). Then
    \[
        f(g_1 \circ g_2) = f(g_1)f(g_2) = e_H,
    \]
    so \(g_1 \circ g_2 \in \ker f\). Also,
    \[
        f(g^{-1}) = f(g)^{-1} = e_H
    \]
    for every \(g \in \ker f\), hence \(g^{-1} \in \ker f\). Therefore \(\ker f\) is a subgroup.

    Now let \(h_1,h_2 \in \operatorname{im} f\). Then there exist \(g_1,g_2 \in G\) such that
    \(f(g_1)=h_1\) and \(f(g_2)=h_2\). Hence
    \[
        h_1h_2 = f(g_1)f(g_2) = f(g_1 \circ g_2) \in \operatorname{im} f.
    \]
    Moreover, if \(h = f(g) \in \operatorname{im} f\), then
    \[
        h^{-1} = f(g)^{-1} = f(g^{-1}) \in \operatorname{im} f.
    \]
    Thus \(\operatorname{im} f\) is a subgroup.
}

\ex{
    Let \(\Lambda \subseteq \mathbb{R}^n\) be a lattice. Then there is a group
    homomorphism
    \[
        \operatorname{Sym}(\Lambda) \to \mathrm{O}(n),
        \qquad
        f(x)=Ax+b \mapsto A.
    \]
    The kernel is the group of translational symmetries of \(\Lambda\), and the
    image is called the \textbf{point group}.
}

\ex{
    There is a group homomorphism
    \[
        \operatorname{sign} : S_n \to \{\pm 1\}.
    \]
    It records whether a permutation is even or odd: even permutations are sent
    to \(1\), odd permutations to \(-1\).

    For example,
    \[
        \sigma = \permu{1,2,3,4}{2,4,1,3}
        \qquad \Rightarrow \qquad
        \operatorname{sign}(\sigma)=-1,
    \]
    while
    \[
        \tilde{\sigma} = \permu{1,2,3,4}{2,3,1,4}
        \qquad \Rightarrow \qquad
        \operatorname{sign}(\tilde{\sigma})=1.
    \]
}

\defn{Alternating group}{
    The \textbf{alternating group} \(A_n\) is the kernel of the sign homomorphism:
    \[
        A_n := \ker(\operatorname{sign}) \leq S_n.
    \]
    Equivalently, \(A_n\) is the subgroup of even permutations.
}

\ex{
    For \(n=3\),
    \[
        A_3
        =
        \left\{
            \permu{1,2,3}{1,2,3},
            \permu{1,2,3}{2,3,1},
            \permu{1,2,3}{3,1,2}
        \right\}.
    \]
    Hence \(A_3\) has three elements and is isomorphic to \(\mathbb{Z}/3\mathbb{Z}\).
}

\section{Conjugacy Classes}

\defn{Conjugation}{
    Let \(G\) be a group and \(g \in G\). Define the group homomorphism
    \[
        \operatorname{Ad}_g : G \to G,
        \qquad
        h \mapsto ghg^{-1}.
    \]
    This is called \textbf{conjugation by \(g\)}.

    Indeed,
    \[
        \operatorname{Ad}_g(e) = geg^{-1} = e,
    \]
    and
    \[
        \operatorname{Ad}_g(hh')
        = ghh'g^{-1}
        = ghg^{-1}gh'g^{-1}
        = \operatorname{Ad}_g(h)\operatorname{Ad}_g(h').
    \]
}

\defn{Conjugate Elements}{
    Let \(h,h' \in G\). We say that \(h\) and \(h'\) are \textbf{conjugate} if
    there exists \(g \in G\) such that
    \[
        h = gh'g^{-1} = \operatorname{Ad}_g(h').
    \]
}

\propp{
    Being conjugate is an equivalence relation on \(G\).
}{
    \begin{itemize}
        \item Reflexivity: every element is conjugate to itself, since
        \(h = ehe^{-1}\).
        \item Symmetry: if \(h \sim h'\), so \(h = gh'g^{-1}\), then
        \(h' = g^{-1}hg\), hence \(h' \sim h\).
        \item Transitivity: if \(h_1 \sim h_2\) and \(h_2 \sim h_3\), say
        \(h_1 = g_1 h_2 g_1^{-1}\) and \(h_2 = g_2 h_3 g_2^{-1}\), then
        \[
            h_1 = g_1 g_2 h_3 (g_1 g_2)^{-1},
        \]
        so \(h_1 \sim h_3\).
    \end{itemize}
}

\defn{Conjugacy Classes}{
    The equivalence classes of the conjugation relation are called
    \textbf{conjugacy classes}. The set of all conjugacy classes is denoted
    by \(\operatorname{conj}(G)\).
}

\ex{
    The conjugacy classes of \(S_3\) are
    \[
        \begin{aligned}
            \operatorname{conj}(S_3)
            ={}& \left\{
                \{e\},
                \left\{
                    \permu{1,2,3}{2,1,3},
                    \permu{1,2,3}{3,2,1},
                    \permu{1,2,3}{1,3,2}
                \right\},
            \right. \\
            & \left.
                \left\{
                    \permu{1,2,3}{2,3,1},
                    \permu{1,2,3}{3,1,2}
                \right\}
            \right\}.
        \end{aligned}
    \]
}

\rmkb{
    A useful way to think about conjugacy classes is that conjugation changes the
    \emph{description} of an object, but not the underlying structure.

    For matrices, if \(M \in \mathrm{GL}(n,\mathbb{R})\), then
    \[
        U M U^{-1}
    \]
    is the same linear map written in a different basis.

    For permutations, if \(\sigma \in S_n\), then
    \[
        \tau \sigma \tau^{-1}
    \]
    simply relabels the elements of \(\{1,\dots,n\}\). This is why the conjugacy
    classes in \(S_n\) are determined by cycle type.

    Normal subgroups are exactly those subgroups that are stable under
    conjugation: if \(N \trianglelefteq G\), then \(gng^{-1} \in N\) for all
    \(g \in G\) and \(n \in N\). This stability is what makes the quotient
    \(G/N\) well defined.
}

\defn{Abelian Group}{
    A group \(G\) is called \textbf{abelian} if
    \[
        hg = gh
        \qquad \forall h,g \in G.
    \]
}

\cor{
    \[
        \operatorname{conj}(G) = G
        \qquad \Longleftrightarrow \qquad
        G \text{ is abelian.}
    \]
}

\defn{Cycle}{
    For \(1 \leq k \leq n\), let \(a_1,\dots,a_k \in \{1,2,\dots,n\}\) be pairwise
    distinct. A \textbf{cycle} \((a_1,\dots,a_k)\) is the permutation
    \[
        a_1 \mapsto a_2,\quad a_2 \mapsto a_3,\quad \dots,\quad a_k \mapsto a_1,
    \]
    which acts as the identity on all other elements. We call \(k\) the
    \textbf{length} of the cycle.
}

\ex{
    To describe \(\operatorname{conj}(S_n)\), it is useful to use cycle type.

    For \(\sigma \in S_n\), the \textbf{cycle type} of \(\sigma\) is the list of
    lengths of the cycles in a disjoint cycle decomposition of \(\sigma\), written
    in decreasing order. Thus the cycle type is a partition
    \[
        \lambda_1 \geq \cdots \geq \lambda_k > 0,
        \qquad
        \sum_{i=1}^k \lambda_i = n.
    \]

    For example, \((1\,3\,2)(4\,5) \in S_5\) has cycle type \((3,2)\).
}

\thmp{Conjugacy In \(S_n\)}{
    Two permutations \(\sigma,\sigma' \in S_n\) are conjugate if and only if they
    have the same cycle type.
}{
    Let \(\tau,\sigma \in S_n\). Then
    \[
        (\tau \sigma \tau^{-1})(\tau(i)) = \tau(\sigma(i)).
    \]
    So conjugation by \(\tau\) simply applies \(\tau\) to the entries in each cycle.
    In particular,
    \[
        \tau (a_1\,\dots\,a_r)\tau^{-1}
        =
        (\tau(a_1)\,\dots\,\tau(a_r)).
    \]
    Hence conjugate permutations have the same cycle type.

    Conversely, if \(\sigma\) and \(\sigma'\) have the same cycle type, then their
    disjoint cycles can be matched cycle by cycle. Choosing \(\tau\) so that it
    sends the entries of each cycle of \(\sigma\) to the corresponding entries of
    the matching cycle of \(\sigma'\), we get
    \[
        \tau \sigma \tau^{-1} = \sigma'.
    \]
}

\ex{
    The conjugacy classes of \(S_4\) are in bijection with the partitions of \(4\):
    \[
        (1,1,1,1), \qquad (2,1,1), \qquad (3,1), \qquad (2,2), \qquad (4).
    \]
    Representatives are:
    \begin{itemize}
        \item \((1,1,1,1)\): \(e\),
        \item \((2,1,1)\): a transposition, e.g. \((1\,2)\),
        \item \((3,1)\): a \(3\)-cycle, e.g. \((1\,2\,3)\),
        \item \((2,2)\): a product of two disjoint transpositions, e.g. \((1\,2)(3\,4)\),
        \item \((4)\): a \(4\)-cycle, e.g. \((1\,2\,3\,4)\).
    \end{itemize}
}

\section{Normal Subgroups}

\defn{Normal Subgroup}{
    A subgroup \(N \leq G\) is called \textbf{normal} if
    \[
        gng^{-1} \in N
        \qquad \forall n \in N,\ \forall g \in G.
    \]
}

\ex{
    In \(S_3\), the subgroup
    \[
        \overline{N} = \{ e, (1\,2) \}
    \]
    is not normal, since
    \[
        (1\,3)(1\,2)(1\,3)^{-1} = (2\,3) \notin \overline{N}.
    \]

    On the other hand,
    \[
        N = \{ e, (1\,2\,3), (1\,3\,2) \}
    \]
    is a normal subgroup of \(S_3\).
}

\thmp{Quotient Group}{
    Let \(N \leq G\) be a normal subgroup. Let \(G/N\) be the set of equivalence
    classes for the relation \(g \sim gn\) for \(n \in N\). Then \(G/N\) is a group
    with multiplication
    \[
        [g]\cdot[h] = [gh].
    \]
}{
    We first check that the multiplication is well defined. If \(g' = gn_1\) and
    \(h' = hn_2\) with \(n_1,n_2 \in N\), then
    \[
        g'h' = gn_1hn_2 = gh\,(h^{-1}n_1h)\,n_2.
    \]
    Since \(N\) is normal, \(h^{-1}n_1h \in N\), hence
    \((h^{-1}n_1h)n_2 \in N\). Therefore \(g'h' \sim gh\), so
    \([g'][h'] = [gh]\).

    The group axioms are inherited from \(G\): associativity follows from that in
    \(G\), the neutral element is \([e]\), and the inverse of \([g]\) is \([g^{-1}]\).

    The classes \([g] = gN \subseteq G\) are also called \textbf{cosets}.
}

\ex{
    Let \(\Lambda \subseteq \mathbb{R}^n\) be a lattice. Then the group of
    translational symmetries of \(\Lambda\) is normal, and
    \[
        \operatorname{Sym}(\Lambda) / \operatorname{Trans}(\Lambda)
    \]
    is the point group.

    Another important example is \(\mathbb{Z}^2 \leq \mathbb{R}^2\). Then
    \(\mathbb{Z}^2\) is a normal subgroup of \((\mathbb{R}^2,+)\), and the quotient
    \[
        T^2 = \mathbb{R}^2 / \mathbb{Z}^2
    \]
    is the \(2\)-torus.
}

\thmp{Isomorphism Theorem}{
    Let \(\varphi : G \to H\) be a group homomorphism.
    \begin{enumerate}
        \item The kernel \(\ker \varphi\) is a normal subgroup of \(G\).
        \item The map
        \[
            \Psi : G/\ker \varphi \to \operatorname{im}\varphi,
            \qquad
            [g] \mapsto \varphi(g),
        \]
        is a well-defined group isomorphism.
    \end{enumerate}
}{
    Let \(n \in \ker \varphi\) and \(g \in G\). Then
    \[
        \varphi(gng^{-1})
        = \varphi(g)\varphi(n)\varphi(g)^{-1}
        = \varphi(g)e_H\varphi(g)^{-1}
        = e_H,
    \]
    so \(gng^{-1} \in \ker \varphi\). Hence \(\ker \varphi \trianglelefteq G\).

    Now define
    \[
        \Psi([g]) := \varphi(g).
    \]
    To see that \(\Psi\) is well defined, let \(g' = gn\) with \(n \in \ker \varphi\). Then
    \[
        \varphi(g') = \varphi(gn) = \varphi(g)\varphi(n) = \varphi(g).
    \]
    Also,
    \[
        \Psi([g][h]) = \Psi([gh]) = \varphi(gh) = \varphi(g)\varphi(h)
        = \Psi([g])\Psi([h]),
    \]
    so \(\Psi\) is a group homomorphism.

    It is surjective by definition of \(\operatorname{im}\varphi\). Finally, if
    \(\Psi([g]) = \Psi([g'])\), then
    \[
        \varphi(g) = \varphi(g')
        \qquad \Longrightarrow \qquad
        \varphi(g^{-1}g') = e_H,
    \]
    hence \(g^{-1}g' \in \ker \varphi\), so \(g' = g(g^{-1}g')\) and therefore
    \([g] = [g']\). Thus \(\Psi\) is injective.
}

\ex{
    Examples:
    \begin{itemize}
        \item If \(f : G \to H\) is surjective, then \(\operatorname{im} f = H\), so
        the isomorphism theorem gives
        \[
            G / \ker f \cong H.
        \]

        \item Every finite abelian group \(A\) is a quotient of \(\mathbb{Z}^n\) for
        some \(n\). Hence
        \[
            A \cong \mathbb{Z}^n / \ker f
        \]
        for a suitable surjective homomorphism \(f : \mathbb{Z}^n \to A\).

        \item Let \(\Lambda \subseteq \mathbb{R}^n\) be a lattice. The homomorphism
        \[
            \operatorname{Sym}(\Lambda) \to \mathrm{O}(n),
            \qquad
            f(x) = Ax+b \mapsto A,
        \]
        has kernel \(\operatorname{Trans}(\Lambda)\). Therefore
        \[
            \operatorname{Sym}(\Lambda) / \operatorname{Trans}(\Lambda)
            \cong
            \operatorname{im} f,
        \]
        and \(\operatorname{im} f\) is the point group of \(\Lambda\).

        \item The determinant defines a surjective homomorphism
        \[
            \det : \mathrm{GL}(n,\mathbb{R}) \to \mathbb{R}_{>0},
        \]
        where \(\mathbb{R}_{>0}\) is viewed as a group under multiplication, and
        \[
            \ker(\det) = \mathrm{SL}(n,\mathbb{R}).
        \]
        Hence
        \[
            \mathrm{GL}(n,\mathbb{R}) / \mathrm{SL}(n,\mathbb{R})
            \cong
            \mathbb{R}_{>0}.
        \]
    \end{itemize}
}

\section{Extensions}

\defn{Exact Sequence}{
    A sequence of groups and group homomorphisms
    \[
        \cdots \to G_i \xrightarrow{f_i} G_{i+1} \xrightarrow{f_{i+1}} G_{i+2} \to \cdots
    \]
    is called \textbf{exact} if
    \[
        \operatorname{im}(f_i) = \ker(f_{i+1})
    \]
    at every term.
}

\rmkb{
    If \(N \trianglelefteq G\), then the quotient map
    \[
        \pi : G \to G/N,
        \qquad
        g \mapsto [g],
    \]
    fits into the exact sequence
    \[
        1 \to N \xrightarrow{\iota} G \xrightarrow{\pi} G/N \to 1,
    \]
    where \(\iota : N \hookrightarrow G\) is the inclusion. Indeed,
    \(\operatorname{im}(\iota) = N = \ker(\pi)\).

    By the isomorphism theorem, any short exact sequence
    \[
        1 \to N \xrightarrow{\iota} G \xrightarrow{\pi} H \to 1
    \]
    identifies \(N\) with a normal subgroup of \(G\) and \(H\) with the quotient
    \(G/N\).
}

\defn{Extension}{
    An \textbf{extension} of a group \(H\) by a group \(N\) is a short exact sequence
    \[
        1 \to N \xrightarrow{\iota} G \xrightarrow{\pi} H \to 1.
    \]
}

\defn{Split And Central Extensions}{
    An extension
    \[
        1 \to N \xrightarrow{\iota} G \xrightarrow{\pi} H \to 1
    \]
    is called \textbf{split} if there exists a group homomorphism called a \textbf{section}
    \[
        s : H \to G
    \]
    such that
    \[
        \pi \circ s = \operatorname{id}_H.
    \]

    It is called \textbf{central} if
    \[
        \iota(n)g = g\iota(n)
        \qquad \forall n \in N,\ \forall g \in G,
    \]
    i.e. if \(\iota(N)\) lies in the center \(Z(G)\) of \(G\).
}

\rmkb{
    A split extension is the group-theoretic analogue of a trivial bundle. Indeed,
    the map \(\pi : G \to H\) may be viewed as a bundle over \(H\) with fiber \(N\),
    and a splitting \(s : H \to G\) with \(\pi \circ s = \operatorname{id}_H\) is a
    global section. The existence of such a distinguished section means that the
    extension can be reconstructed from \(N\) and \(H\) without any twisting; in
    the simplest case one gets the direct product \(N \times H\), and more generally
    a semidirect product.
}

\rmkb{
    Central extensions are especially important in quantization. In classical
    mechanics one often starts with a symmetry group acting on the phase space,
    but after quantization the corresponding operators on the Hilbert space may
    realize the symmetry only up to phases. These extra phase factors are encoded
    by a central extension, typically by \(U(1)\) or \(\mathbb{R}\). In this sense,
    central extensions measure the obstruction to lifting a classical symmetry
    action to an honest linear action in quantum theory.
}

\ex{
    Examples:
    \begin{itemize}
        \item The product group \(N \times H\) gives an extension
        \[
            1 \to N \to N \times H \to H \to 1,
        \]
        which is always split. If \(N\) is abelian, then this extension is also central.

        \item The symmetry group of a lattice fits into an extension
        \[
            1 \to \operatorname{Trans}(\Lambda) \to \operatorname{Sym}(\Lambda)
            \to \text{point group} \to 1.
        \]

        \item The Lorentz group fits into the extension
        \[
            1 \to \mathrm{SO}^+(3,1) \to \mathrm{O}(3,1) \to \mathbb{Z}_2 \times \mathbb{Z}_2 \to 1,
        \]
        which is split.

        \item The Poincare group \(P\) fits into the extension
        \[
            1 \to \mathbb{R}^4 \to P \to \mathrm{O}(3,1) \to 1,
        \]
        which is split.

        \item The sequence
        \[
            1 \to \mathbb{Z}_2 \to \mathbb{Z}_4 \to \mathbb{Z}_2 \to 1
        \]
        is a central extension, but it is not split.
    \end{itemize}
}

\defn{Semidirect Product}{
    Let \(\psi : H \to \operatorname{Aut}(N)\) be a group homomorphism. The
    \textbf{semidirect product} \(N \rtimes_{\psi} H\) is the group whose underlying
    set is \(N \times H\), with multiplication
    \[
        (n,h)\cdot(n',h')
        :=
        \bigl(n\,\psi(h)[n'],\,hh'\bigr).
    \]

    In this notation one may think of the relation
    \[
        hn = \psi(h)[n]\,h.
    \]
}

\rmkb{
    There is a canonical split extension
    \[
        1 \to N \xrightarrow{\iota} N \rtimes_{\psi} H \xrightarrow{\pi} H \to 1,
    \]
    where
    \[
        \iota(n) = (n,e_H),
        \qquad
        \pi(n,h) = h,
        \qquad
        s(h) = (e_N,h).
    \]
    Thus every semidirect product gives a split extension, and conversely every
    split extension is isomorphic to a suitable semidirect product.
}

\thmp{Split Extensions And Semidirect Products}{
    Let
    \[
        1 \to N \xrightarrow{\iota} G \xrightarrow{\pi} H \to 1
    \]
    be an extension, and let \(s : H \to G\) be a splitting. Then:
    \begin{enumerate}
        \item Conjugation by \(s(h)\) defines a group homomorphism
        \[
            \psi : H \to \operatorname{Aut}(N),
            \qquad
            h \mapsto \bigl(n \mapsto \iota^{-1}(s(h)\iota(n)s(h)^{-1})\bigr).
        \]
        \item The group \(G\) is isomorphic to the semidirect product
        \[
            N \rtimes_{\psi} H.
        \]
    \end{enumerate}
}{
    Since \(\iota(N)\trianglelefteq G\), conjugation by \(s(h)\) preserves
    \(\iota(N)\), hence defines an automorphism of \(N\). This gives the map
    \(\psi : H \to \operatorname{Aut}(N)\), and it is a group homomorphism because
    conjugation is compatible with multiplication.

    Define
    \[
        \Phi : N \rtimes_{\psi} H \to G,
        \qquad
        (n,h) \mapsto \iota(n)s(h).
    \]
    Using the definition of \(\psi\), one checks
    \[
        \Phi\bigl((n,h)(n',h')\bigr)
        = \Phi(n,h)\Phi(n',h').
    \]
    Thus \(\Phi\) is a homomorphism.

    It is surjective since every \(g \in G\) can be written as
    \(g = \iota(n)s(h)\), where \(h = \pi(g)\) and \(n\) is determined by the fact
    that \(g s(h)^{-1} \in \ker \pi = \iota(N)\). It is injective because
    \(\iota(n)s(h) = e\) implies \(h = e\) after applying \(\pi\), and then \(n=e\).
    Hence \(\Phi\) is an isomorphism.
}

\rmkb{
    If an extension is both split and central, then the associated action
    \(\psi : H \to \operatorname{Aut}(N)\) is trivial, because
    \[
        \iota(n)s(h) = s(h)\iota(n)
        \qquad \forall n \in N,\ h \in H.
    \]
    Therefore
    \[
        N \rtimes_{\psi} H = N \times H,
    \]
    so every split central extension is in fact a direct product.
}





